\documentclass{article}
\usepackage{amsmath,amssymb,amsthm}
\usepackage{algorithm}
\usepackage{algpseudocode}
\usepackage{fullpage}
\usepackage{xcolor}
\newcommand{\E}{\mathbb{E}}
\newcommand{\R}{\mathbb{R}}
\newcommand{\norm}[1]{\left\| #1 \right\|}
\newcommand{\inner}[1]{\left\langle #1 \right\rangle}
\newtheorem{assumption}{Assumption}
\newtheorem{theorem}{Theorem}
\newtheorem{lemma}{Lemma}
\newtheorem{corollary}{Corollary}
\begin{document}

\section*{Cyclic Stochastic Gradient Descent (Cyclic SGD)}

\section*{Mathematical Formulation}
Let $f:\R^{d}\to\R$ be the (possibly non‑convex) objective function and let $\{z_t\}_{t\ge0}$ be an i.i.d. data stream.  
At iteration $t$ we compute a stochastic gradient
\[
g_t \;:=\; \nabla f(w_t;z_t)\in\R^{d},
\]
and update the parameter vector $w_t\in\R^{d}$ by
\begin{equation}
\boxed{%
w_{t+1}=w_t-\eta_t\, g_t
}\tag{1}
\end{equation}
where the **cyclic learning‑rate schedule** $\{\eta_t\}$ is defined as follows.
Given three hyper‑parameters
\[
0<\eta_{\min}<\eta_{\max},\qquad M\in\mathbb{N},
\]
let $s_t:=t\;\bmod\; 2M$ and set
\[
\eta_t=
\begin{cases}
\displaystyle 
\eta_{\min}+(\eta_{\max}-\eta_{\min})\frac{s_t}{M},
& 0\le s_t<M,\\[2ex]
\displaystyle 
\eta_{\max}-(\eta_{\max}-\eta_{\min})\frac{s_t-M}{M},
& M\le s_t<2M .
\end{cases}
\tag{2}
\]
Thus $\eta_t$ increases linearly from $\eta_{\min}$ to $\eta_{\max}$ in $M$ steps and then decreases back to $\eta_{\min}$ in the next $M$ steps, forming a triangular wave that repeats every $2M$ iterations.

\section*{Pseudocode}
\begin{algorithm}[H]
\caption{Cyclic Stochastic Gradient Descent}
\begin{algorithmic}[1]
\State \textbf{Input:} initial point $w_0\in\R^{d}$, \; $\eta_{\min},\eta_{\max}>0$, cycle length $M\in\mathbb{N}$, \; max iterations $T$.
\For{$t=0,1,\dots,T-1$}
    \State Compute $s_t \gets t \bmod 2M$.
    \If{$s_t < M$}
        \State $\eta_t \gets \eta_{\min}+(\eta_{\max}-\eta_{\min})\frac{s_t}{M}$ \Comment{ascending phase}
    \Else
        \State $\eta_t \gets \eta_{\max}-(\eta_{\max}-\eta_{\min})\frac{s_t-M}{M}$ \Comment{descending phase}
    \EndIf
    \State Sample a data point $z_t$ (or minibatch) and evaluate stochastic gradient $g_t\gets \nabla f(w_t;z_t)$.
    \State Update $w_{t+1}\gets w_t-\eta_t\,g_t$.
\EndFor
\State \textbf{Output:} $w_T$ (or optionally the averaged iterate $\bar w_T=\frac1T\sum_{t=0}^{T-1}w_t$).
\end{algorithmic}
\end{algorithm}

\section*{Dry Run Example}
Consider the scalar quadratic $f(w)=(w-3)^2$ with exact gradients (no stochasticity).  
Set $w_0=0$, $\eta_{\min}=0.05$, $\eta_{\max}=0.15$, $M=2$; thus the schedule for the first four steps is
\[
\eta_0=0.05,\quad \eta_1=0.10,\quad \eta_2=0.15,\quad \eta_3=0.10.
\]

\begin{center}
\begin{tabular}{c|c|c|c}
$t$ & $\eta_t$ & $g_t=2(w_t-3)$ & $w_{t+1}=w_t-\eta_t g_t$\\\hline
0 & $0.05$ & $2(0-3)=-6$ & $w_1=0-0.05(-6)=0.30$\\
1 & $0.10$ & $2(0.30-3)=-5.40$ & $w_2=0.30-0.10(-5.40)=0.84$\\
2 & $0.15$ & $2(0.84-3)=-4.32$ & $w_3=0.84-0.15(-4.32)=1.50$
\end{tabular}
\end{center}
The corresponding objective values are
\[
f(w_0)=9,\quad f(w_1)=7.29,\quad f(w_2)=5.86,\quad f(w_3)=4.50,
\]
illustrating a monotone decrease despite the cyclic step‑size.

\newpage

\section*{Assumptions}
\begin{assumption}[Smoothness]\label{ass:smooth}
$f$ is $L$‑smooth, i.e., for all $x,y\in\R^{d}$,
\[
f(y)\le f(x)+\inner{\nabla f(x)}{y-x}+\frac{L}{2}\norm{y-x}^2 .
\]
\end{assumption}
\begin{assumption}[Unbiased Gradient]\label{ass:unbiased}
For every $t$, conditional on $w_t$,
\[
\E[g_t\mid w_t]=\nabla f(w_t).
\]
\end{assumption}
\begin{assumption}[Bounded Variance]\label{ass:variance}
There exists $\sigma^2\ge0$ such that
\[
\E\big[\norm{g_t-\nabla f(w_t)}^2\mid w_t\big]\le\sigma^2,\qquad\forall t .
\]
\end{assumption}
\begin{assumption}[Step‑size Bounds]\label{ass:stepsize}
The cyclic schedule satisfies
\[
0<\eta_{\min}\le\eta_t\le\eta_{\max}\le\frac{1}{L},\qquad\forall t .
\]
\end{assumption}

\section*{Main Convergence Theorem}
\begin{theorem}[Convergence of Cyclic SGD]\label{thm:main}
Under Assumptions~\ref{ass:smooth}–\ref{ass:stepsize}, let $\{w_t\}_{t=0}^{T}$ be generated by \eqref{1}–\eqref{2}. Define the average step size
\[
\bar\eta_T:=\frac1T\sum_{t=0}^{T-1}\eta_t\;\;\in\;[\eta_{\min},\eta_{\max}].
\]
Then the expected squared norm of the gradient satisfies
\begin{equation}
\frac1T\sum_{t=0}^{T-1}\E\!\big[\norm{\nabla f(w_t)}^2\big]
\;\le\;
\frac{2\big(f(w_0)-f^{\star}\big)}{\eta_{\min}T}
\;+\;
L\,\eta_{\max}\,\sigma^2,
\tag{3}
\end{equation}
where $f^{\star}:=\inf_{w}f(w)$. Consequently,
\[
\liminf_{T\to\infty}\frac1T\sum_{t=0}^{T-1}\E\!\big[\norm{\nabla f(w_t)}^2\big]\le L\,\eta_{\max}\,\sigma^2 .
\]
\end{theorem}

\section*{Theoretical Properties}
\begin{itemize}
\item \textbf{Stability:} The upper bound \eqref{3} is finite as long as $\eta_{\max}\le 1/L$, guaranteeing that the iterates cannot diverge due to excessively large steps.
\item \textbf{Hyper‑parameter Sensitivity:} The term $L\eta_{\max}\sigma^2$ reflects the steady‑state error; decreasing $\eta_{\max}$ (or the variance $\sigma^2$) reduces the asymptotic gradient norm.
\item \textbf{Comparison to Constant‑Step SGD:} For a constant step $\eta$, the bound becomes $\frac{2(f(w_0)-f^{\star})}{\eta T}+L\eta\sigma^2$. The cyclic schedule replaces the single $\eta$ by $\eta_{\min}$ in the first term (worst‑case) while the second term depends on the maximal step $\eta_{\max}$. Thus cyclic SGD inherits the same $O(1/T)$ decay as constant‑step SGD but may enjoy better practical exploration because of the periodic increase to $\eta_{\max}$.
\item \textbf{Special Cases:}
    \begin{enumerate}
        \item If $\sigma^2=0$ (full‑batch gradients), the bound reduces to $\frac{2(f(w_0)-f^{\star})}{\eta_{\min}T}$, yielding exact convergence to a stationary point.
        \item If $\eta_{\min}=\eta_{\max}=\eta$, the theorem collapses to the classical constant‑step SGD result.
    \end{enumerate}
\end{itemize}

\section*{Discussion}
The theorem shows that even though the learning rate is *non‑monotonic* and revisits large values, the average progress per iteration is governed by the smallest step size $\eta_{\min}$ (appearing in the $1/T$ term) while the variance‑induced error is controlled by the largest step $\eta_{\max}$. This dual dependence explains why practitioners often observe that cyclic schedules accelerate early learning (large $\eta_{\max}$) without sacrificing asymptotic stability (small $\eta_{\min}$).

\newpage

\section*{Preliminaries (Definitions and Lemmas)}
\begin{lemma}[One‑step Smoothness Inequality]\label{lem:smooth}
If $f$ is $L$‑smooth, then for any $w$ and any vector $u$,
\[
f(w-u)\le f(w)-\inner{\nabla f(w)}{u}+\frac{L}{2}\norm{u}^2 .
\]
\end{lemma}
\begin{proof}
Apply Assumption~\ref{ass:smooth} with $x=w$, $y=w-u$.
\end{proof}

\begin{lemma}[Variance Decomposition]\label{lem:var}
For any random vector $g$ with $\E[g]=\mu$,
\[
\E[\norm{g}^2]=\norm{\mu}^2+\E[\norm{g-\mu}^2].
\]
\end{lemma}
\begin{proof}
Expand $\norm{g}^2=\norm{\mu+(g-\mu)}^2$ and take expectation.
\end{proof}

\section*{Main Proof (Step‑by‑Step Derivation)}
We prove Theorem~\ref{thm:main}.  All expectations $\E[\cdot]$ are taken over the randomness of the stochastic gradients up to the current iteration unless otherwise noted.

\noindent\textbf{[Step 1]} Apply Lemma~\ref{lem:smooth} with $u=\eta_t g_t$:
\[
f(w_{t+1})\le f(w_t)-\eta_t\inner{\nabla f(w_t)}{g_t}
+\frac{L}{2}\eta_t^{2}\norm{g_t}^{2}.
\tag{4}
\]

\noindent\textbf{[Step 2]} Take conditional expectation w.r.t.\ $w_t$ on both sides of \eqref{4} and use Assumption~\ref{ass:unbiased}:
\[
\E\!\big[f(w_{t+1})\mid w_t\big]\le
f(w_t)-\eta_t\norm{\nabla f(w_t)}^{2}
+\frac{L}{2}\eta_t^{2}\E\!\big[\norm{g_t}^{2}\mid w_t\big].
\tag{5}
\]

\noindent\textbf{[Step 3]} Apply Lemma~\ref{lem:var} to $\E[\norm{g_t}^{2}\mid w_t]$:
\[
\E\!\big[\norm{g_t}^{2}\mid w_t\big]
= \norm{\nabla f(w_t)}^{2}
+ \E\!\big[\norm{g_t-\nabla f(w_t)}^{2}\mid w_t\big]
\le \norm{\nabla f(w_t)}^{2}+\sigma^{2},
\tag{6}
\]
where the inequality follows from Assumption~\ref{ass:variance}.

\noindent\textbf{[Step 4]} Substitute \eqref{6} into \eqref{5}:
\[
\E\!\big[f(w_{t+1})\mid w_t\big]\le
f(w_t)-\eta_t\norm{\nabla f(w_t)}^{2}
+\frac{L}{2}\eta_t^{2}\big(\norm{\nabla f(w_t)}^{2}+\sigma^{2}\big).
\tag{7}
\]

\noindent\textbf{[Step 5]} Rearrange \eqref{7} to isolate the gradient term:
\[
\Big(\eta_t-\frac{L}{2}\eta_t^{2}\Big)\norm{\nabla f(w_t)}^{2}
\le f(w_t)-\E\!\big[f(w_{t+1})\mid w_t\big]
+\frac{L}{2}\eta_t^{2}\sigma^{2}.
\tag{8}
\]

\noindent\textbf{[Step 6]} By Assumption~\ref{ass:stepsize}, $\eta_t\le 1/L$, hence
\[
\eta_t-\frac{L}{2}\eta_t^{2}
\;\ge\; \frac{\eta_t}{2}
\;\ge\; \frac{\eta_{\min}}{2}.
\tag{9}
\]
(The first inequality uses $1-\frac{L}{2}\eta_t\ge\frac12$ because $L\eta_t\le1$.)

\noindent\textbf{[Step 7]} Using \eqref{9} in \eqref{8} gives
\[
\frac{\eta_{\min}}{2}\,\norm{\nabla f(w_t)}^{2}
\le f(w_t)-\E\!\big[f(w_{t+1})\mid w_t\big]
+\frac{L}{2}\eta_{\max}^{2}\sigma^{2}.
\tag{10}
\]

\noindent\textbf{[Step 8]} Take full expectation (over all randomness up to time $t$) on \eqref{10}:
\[
\frac{\eta_{\min}}{2}\,\E\!\big[\norm{\nabla f(w_t)}^{2}\big]
\le \E\!\big[f(w_t)\big]-\E\!\big[f(w_{t+1})\big]
+\frac{L}{2}\eta_{\max}^{2}\sigma^{2}.
\tag{11}
\]

\noindent\textbf{[Step 9]} Sum \eqref{11} over $t=0,\dots,T-1$:
\[
\frac{\eta_{\min}}{2}\sum_{t=0}^{T-1}\E\!\big[\norm{\nabla f(w_t)}^{2}\big]
\le f(w_0)-\E\!\big[f(w_T)\big]
+\frac{L}{2}\eta_{\max}^{2}\sigma^{2}T .
\tag{12}
\]

\noindent\textbf{[Step 10]} Since $f(w_T)\ge f^{\star}$, replace $\E[f(w_T)]$ by $f^{\star}$ to obtain an upper bound:
\[
\frac{\eta_{\min}}{2}\sum_{t=0}^{T-1}\E\!\big[\norm{\nabla f(w_t)}^{2}\big]
\le f(w_0)-f^{\star}
+\frac{L}{2}\eta_{\max}^{2}\sigma^{2}T .
\tag{13}
\]

\noindent\textbf{[Step 11]} Divide both sides of \eqref{13} by $\eta_{\min}T$ and rearrange:
\[
\frac1T\sum_{t=0}^{T-1}\E\!\big[\norm{\nabla f(w_t)}^{2}\big]
\le
\frac{2\big(f(w_0)-f^{\star}\big)}{\eta_{\min}T}
+L\,\eta_{\max}\sigma^{2}.
\tag{14}
\]
Equation \eqref{14} is exactly the bound \eqref{3} claimed in Theorem~\ref{thm:main}.  $\square$

\section*{Rate Derivation (With Constants)}
The bound \eqref{3} consists of two terms:
\begin{itemize}
\item A \emph{decaying term} $\displaystyle \frac{2(f(w_0)-f^{\star})}{\eta_{\min}T}=O\!\big(\frac1T\big)$.
\item A \emph{steady‑state term} $L\eta_{\max}\sigma^{2}=O(\eta_{\max})$.
\end{itemize}
Thus, for a fixed cyclic schedule the algorithm attains an $O(1/T)$ convergence rate to a neighbourhood of stationary points whose radius is proportional to $\sqrt{L\eta_{\max}\sigma^{2}}$.

\section*{Special Cases}
\begin{enumerate}
\item \textbf{Zero variance ($\sigma^{2}=0$).} The steady‑state term vanishes and we recover exact $O(1/T)$ convergence to a stationary point.
\item \textbf{Constant step size.} Setting $\eta_{\min}=\eta_{\max}=\eta$ reduces the bound to the classic SGD result
\[
\frac1T\sum_{t=0}^{T-1}\E\!\big[\norm{\nabla f(w_t)}^{2}\big]
\le\frac{2(f(w_0)-f^{\star})}{\eta T}+L\eta\sigma^{2}.
\]
\item \textbf{Very short cycle ($M=1$).} The schedule becomes constant $\eta_t\equiv\eta_{\min}=\eta_{\max}$, again yielding the constant‑step bound.
\end{enumerate}

\end{document}